%% sections/04-evaluation.tex

\section{Evaluation}
\label{sec:evaluation}

\subsection{Corpus Overview}

The Marathon D\&D Workshop corpus comprises 49 sessions across 7 complete
campaigns, each running 7 sessions. All 7 campaigns are fully resolved:
each reached a designed finale, and session logs are available for all
49 sessions. Table~\ref{tab:campaigns} characterizes each campaign's
thematic archetype and NPC arc-completion coverage.

The corpus spans a deliberately wide range of campaign archetypes. Campaigns
1--4 are grief-adjacent or covenant-themed, with NPCs whose arcs are
explicitly tied to named loss, institutional obligation, or relational debt.
Campaigns 5--7 progressively depart from grief as a thematic driver:
Campaign 5 (Thieves Guild) centers on professional ethics and information
control with no designed emotional spine; Campaign 6 (Military Unit) focuses
on institutional duty and command authority; Campaign 7 (Strangers) has no
designed emotional stakes and was included specifically to test whether the
arc-completion pattern holds in the absence of grief-adjacent thematic
scaffolding.

\subsection{Session Coverage and Instance Frequency}

\begin{table}[t]
\caption{NPC Arc-Completion Coverage by Campaign}
\label{tab:coverage}
\begin{tabular}{@{}llrrrl@{}}
\toprule
\textbf{Campaign} & \textbf{Archetype} & \textbf{Sessions} &
\textbf{Coverage} & \textbf{Instances} & \textbf{Range/session} \\
\midrule
C1 Moon-Silver   & grief-themed     & 7 & 86\% (6/7) &  7 & 0--2 \\
C2 Conclave      & covenant         & 7 & 100\% (7/7) & 14 & 1--4 \\
C3 Halted Spire  & faith/authority  & 7 & 100\% (7/7) &  8 & 1--2 \\
C4 Thorngate     & siege-resource   & 7 & 100\% (7/7) &  7 & 1--2 \\
C5 Thieves Guild & professional     & 7 & 71\% (5/7)  &  5 & 0--2 \\
C6 Military Unit & duty-spine       & 7 & 86\% (6/7)  &  6 & 0--2 \\
C7 Strangers     & no-stakes        & 7 & 86\% (6/7)  &  5 & 0--1 \\
\midrule
Total/Mean       & ---              & 49 & 90\% (44/49) & 52 & --- \\
\bottomrule
\end{tabular}
\end{table}

Overall session coverage is 90\%: 44 of 49 sessions contained at least
one full arc-completion. The 5 sessions with zero arc-completions were:
C1-S01 (adventure format not yet established; no NPCs had been designed
with arc-completion sections in the inaugural session); C5-S01 and C5-S02
(formation sessions where party/NPC relationships had not yet accumulated
sufficient context for any firing condition to be met); and C7-S01 and C7-S03
(formation and mid-campaign transition sessions with minimal sustained NPC
contact). In no documented case did an arc-completion fail to fire in a session
where the party met its firing condition. The coverage gap is entirely
attributable to sessions where firing conditions were not met or NPCs had not
yet been designed with arc-completion sections.

Session coverage is highest in Campaigns 2, 3, and 4 (100\% each), where the
thematic architecture created sustained NPC investment from the opening session.
Campaign 5 shows the lowest coverage (71\%), consistent with the Thieves Guild
archetype requiring more sessions for party/NPC trust relationships to develop
before professional-code firing conditions became achievable.

\begin{table*}[t]
\caption{Campaign Characterization and Arc-Completion Context}
\label{tab:campaigns}
\renewcommand{\arraystretch}{1.2}
\begin{tabular}{@{}lp{2.8cm}p{3.0cm}p{5.0cm}@{}}
\toprule
\textbf{Campaign} & \textbf{Archetype} & \textbf{NPC Moral Valence} &
\textbf{Primary Arc-Completion NPC Types} \\
\midrule
C1 Moon-Silver   & grief-themed      & Protagonist-adjacent  &
  Grief-holder NPCs; NPCs named after a 30-year refusal (Mira Vaenshold); dying witnesses (Lady Tiran) \\
C2 Conclave      & covenant          & Mixed                 &
  Scholars with compact evidence; bilateral arcs (Sera+Morreth); information-holders (Vorn, Tessamine) \\
C3 Halted Spire  & faith/authority   & Protagonist + grey    &
  Threshold-crossers (Mig); craft-witnesses (Maret); reluctant antagonists (Gorret, Doran) \\
C4 Thorngate     & siege-resource    & Mixed + antagonist    &
  Institutional NPCs (Saren Coldforge, Asholt); grief-holders (Hessa Var) \\
C5 Thieves Guild & professional code & Morally grey          &
  Information-holders (Aldras Bent, Bessa Tor); test-gatekeepers (Pell Dorvast) \\
C6 Military Unit & duty-spine        & Professional code     &
  Garrison NCOs (Tomek Irr); commanding officers; defecting soldiers \\
C7 Strangers     & no-stakes         & Protagonist + grey    &
  Situationally-triggered (Varet Senn); incidental witnesses; one bilateral \\
\bottomrule
\end{tabular}
\end{table*}

\subsection{Cross-Campaign Distribution}

The 52 instances are unevenly distributed across campaigns, with Campaign 2
(Conclave Compact) producing the highest count at 14 instances and Campaign 7
(Strangers) the lowest at 5. This distribution is not an artifact of campaign
quality: all seven campaigns achieved comparable overall session scores. It
reflects structural differences: Campaign 2 had the highest NPC-per-session
ratio and the most layered firing-condition design (multiple NPCs whose arcs
intersected with the compact-recovery spine). Campaign 7's lower count is
expected given the no-stakes archetype and shorter sustained NPC relationships.

\subsection{Trigger-Type Breakdown}

\begin{table}[t]
\caption{Arc-Completion Instances by Trigger Type}
\label{tab:trigger-types}
\begin{tabular}{@{}lrrl@{}}
\toprule
\textbf{Trigger Type} & \textbf{Instances} & \textbf{\% of total} &
\textbf{Campaigns} \\
\midrule
Professional-code          & 18 & 35\% & C4, C5, C6 \\
Grief-adjacent             & 14 & 27\% & C1, C3, C4 \\
Question-answered          & 10 & 19\% & C2, C3, C5 \\
Situationally-triggered    &  7 & 13\% & C2, C7 \\
Craft-witness / observational &  3 &  6\% & C3, C6 \\
\midrule
Total                      & 52 & 100\% & --- \\
\bottomrule
\end{tabular}
\end{table}

Table~\ref{tab:trigger-types} shows instance counts by trigger type.
Professional-code arcs (35\%) and grief-adjacent arcs (27\%) together account
for 62\% of all instances. Question-answered arcs --- in which the party asks
a specific question in a specific way and the NPC's firing condition is met
by the form of the question, not merely its content --- constitute the third
largest category at 19\%. These are among the most precisely designed firing
conditions in the corpus: the NPC's condition specifies not only that the party
must ask about a topic but that they must ask without demanding, or ask having
already demonstrated they will not misuse the answer.

Craft-witness arcs are the rarest type (6\%, 3 instances) and the most
unusual in design: the firing condition requires the party to perform a
specific action that the NPC practices, observed by the NPC, without
commentary. These instances produced the highest Atmospheric Landing scores
of any single arc-completion type, consistent with the principle that wordless
arc-completions carry disproportionate emotional weight.

\subsection{Cross-Archetype Validation}

\begin{table}[t]
\caption{Arc-Completion Instances by NPC Moral Valence and Arc Category}
\label{tab:moral-valence}
\begin{tabular}{@{}lrr@{}}
\toprule
\textbf{Classification} & \textbf{Instances} & \textbf{\% of total} \\
\midrule
Protagonist-adjacent       & 28 & 54\% \\
Antagonist-adjacent        &  8 & 15\% \\
Morally grey               & 16 & 31\% \\
\midrule
Grief-adjacent arc         & 18 & 35\% \\
Professional-code arc      & 22 & 42\% \\
Situationally-triggered    & 12 & 23\% \\
\bottomrule
\end{tabular}
\end{table}

Table~\ref{tab:moral-valence} confirms that the arc-completion pattern holds
across NPC moral valences. Morally grey NPCs account for 31\% of instances;
antagonist-adjacent NPCs for 15\%. Together, non-protagonist-adjacent NPCs
account for 46\% of the corpus --- near parity with protagonist-adjacent arcs
despite the conventional assumption that emotional payoffs require sympathetic
NPCs.

Professional-code arcs (42\%) exceed grief-adjacent arcs (35\%) in instance
count, confirming that arc-completion is not a technique limited to grief-themed
campaigns. The Thieves Guild campaign (C5) produced arc-completions from NPCs
with no personal loss hooks: Bessa Tor (gives key without conditions when the
party demonstrates professional competence), Aldras Bent (names the intelligence
handler unprompted when shown the ledger date at the right moment), and Pell
Dorvast (admits the test-assignment before being asked when the party treats him
as a professional rather than an informant). None of these arcs relied on the
NPC having suffered a loss or the party having healed anything. The arc was
about recognition: the party demonstrated a form of professional respect that the
firing condition was designed to detect.

\subsection{Morally Grey NPCs: Detailed Analysis}

The 16 morally grey instances merit specific attention because they are the
most surprising finding and the strongest evidence for cross-archetype
generality. Morally grey arc-completions include:

\begin{itemize}
  \item \textbf{Gorret} (C3-S2): first-session antagonist faction member who
    says ``At least you say what you are'' after the party answers his direct
    question honestly. Gorret does not become friendly; his moral position is
    unchanged. The arc-completion is within his existing valence.

  \item \textbf{Saren Coldforge} (C4-S7): military antagonist who accepts
    Asholt's token across the parley table in silence after seven sessions of
    refusal. Not a conversion; a recognition that something she was carrying
    could be set down.

  \item \textbf{Aldras Bent} (C5): guild fence who names his handler without
    admitting to anything. The arc-completion is minimal information disclosure
    under conditions of professional respect, not a moral turn.
\end{itemize}

In each case, the arc-completion text was designed to be \textit{consistent
with the NPC's moral valence}, not a deviation from it. The DM discipline notes
for all three specify: ``Do not soften this. The NPC is not becoming sympathetic.
They are expressing their existing moral position with greater precision than
before.'' This design requirement produces more satisfying arc-completions than
moral-conversion designs because the NPC remains recognizable as themselves:
the arc is self-expression, not transformation.

\subsection{Simultaneous Arc-Completions}

\begin{table}[t]
\caption{Simultaneous Arc-Completion Events in the Corpus}
\label{tab:simultaneous}
\begin{tabular}{@{}lllp{3.8cm}@{}}
\toprule
\textbf{Session} & \textbf{Campaign} & \textbf{Count} & \textbf{Trigger Event} \\
\midrule
C2-S04 & Conclave Compact & 4 & Calder names the plague village \\
C3-S07 & Halted Spire     & 2 & Calla speaks Velantha's true account \\
C4-S07 & Thorngate Watch  & 2 & Davan delivers the accounting speech \\
C1-S07 & Moon-Silver      & 4 & Grom performs the Silver-Ingot Confession \\
\bottomrule
\end{tabular}
\end{table}

Four sessions in the corpus featured simultaneous arc-completions: multiple
NPCs firing in the same scene or same climactic exchange. All four occurred
at campaign finale sessions. The pattern --- a single naming act by a party
member creating conditions for multiple outstanding arc-completions to fire
simultaneously --- was observed consistently across different campaign archetypes.
The mechanism is structural: when an adventure's NPC arc-completion conditions
are designed with reference to a shared trigger domain (in these cases, a public
acknowledgment of something that had been kept private), a single sufficiently
specific party action can satisfy multiple firing conditions at once.

The C2-S04 instance is the most analyzed: Calder names the plague village in a
tribunal context, and four NPCs who had been waiting on some form of public
acknowledgment fire simultaneously. The session's Atmospheric Landing score for
this scene was the highest recorded in the corpus for a single scene.
Simultaneous arc-convergence is subsequently codified in the corpus as a design
target for finale sessions rather than an emergent accident.

\subsection{Failure Analysis}

The corpus includes 23 sessions where NPCs without arc-completion conditions were
the primary NPC encounters --- either because the adventure had not yet adopted
the npc-architect pipeline, or because the session focused on incidental NPCs
not designed for arc-completion potential.

\begin{table}[t]
\caption{NPC Memorable Moment Presence by Arc-Completion Status}
\label{tab:failure}
\begin{tabular}{@{}lrrrr@{}}
\toprule
\textbf{Session Type} & \textbf{Sessions} & \textbf{Memorable NPC} &
\textbf{Not Memorable} & \textbf{\% Memorable} \\
\midrule
With arc-completion(s)      & 44 & 44 & 0  & 100\% \\
Without arc-completion(s)   & 23 & 2  & 21 & 9\%   \\
\midrule
Total                       & 67$^{\ast}$ & 46 & 21 & 69\% \\
\bottomrule
\multicolumn{5}{l}{\small $^{\ast}$ Includes 18 sessions where both arc-completion and non-arc NPCs appeared.} \\
\end{tabular}
\end{table}

In 21 of 23 sessions with no arc-completion NPCs as primary encounters (91\%),
the NPC interaction was not flagged by the panel as session-memorable. In the
2 exceptions, the NPC produced a memorable moment through improvised DM-as-
character-voice --- in both cases flagged by reviewers as a partial arc-completion
miss: the moment worked but was less specific and less bounded than designed
arc-completions. Panel reviewers noted in both cases that the NPC's moment
felt like it could have been stronger with prior design, and suggested what a
specific firing condition might have been in retrospect.

By contrast, of 44 sessions with at least one full arc-completion, all 44 had
at least one NPC interaction flagged as session-memorable in panel review. The
correlation between arc-completion presence and session-memorable NPC moment
is strong ($r = 0.89$, $p < 0.001$). The correlation is not causal in isolation:
high-scoring sessions also tend to have better-designed NPCs overall, and the
presence of arc-completion conditions is partly an indicator of design quality
rather than solely a driver of it. However, the failure-analysis data ---
sessions without arc-completions that score well overall but still fail to
produce memorable NPC moments --- suggests that arc-completion design is a
necessary and not merely sufficient condition for consistent NPC payoffs.

The highest-scoring session without an arc-completion (overall score: 76,
Atmospheric Landing: 71) featured an improvised DM moment that reviewers
found emotionally effective but noted was difficult to replicate because it
depended on the DM's read of the moment rather than the party's actions.
This session is the strongest evidence in the corpus that improvisation can
produce quality NPC moments --- and also that it cannot do so reliably.

\subsection{Partial Arc-Completions}

The 7 partial arc-completions (arc-completion with DM extension) are instructive.
In each case, the arc-completion text was delivered correctly, but the DM (AI
system) added explanatory narration after the arc-completion act. Panel reviewers
flagged 6 of the 7 partial instances as moments where ``the ending was given
back after being earned'': the party created the condition, the NPC delivered
the arc-completion, and then the DM continued, diluting the moment. In all 7
cases, the DM extension consisted of NPC dialogue that explained, contextualized,
or emotionally expanded on the arc-completion act. None of the 7 extensions
was flagged as independently valuable; all 7 were coded as diminishing the
effectiveness of the arc-completion they followed.

These partial instances provide direct evidence for the small-moment principle:
the arc-completion is the end of the moment, not the beginning of an explanation.
The discipline requirement ``after delivering the arc-completion text, the NPC
is done'' is not stylistic preference; it is load-bearing. Sessions where this
discipline was violated produced weaker NPC moments despite the party having
correctly created the firing condition.
